%\documentclass[10pt,dvipsnames,svgnames]{beamer}
\documentclass[aspectratio=169, lualatex, handout, 10pt,dvipsnames,svgnames]{beamer} %
\makeatletter\def\input@path{{theme/}}\makeatother\usetheme{cipher}

\title{Computer Security: Introduction}
\author{Markulf Kohlweiss}
\subject{Introduction to Applied Cryptography covering cryptographic primitives, protocols, security goals, and real-world applications from basic concepts to advanced topics like TLS and zero-knowledge proofs.}
\keywords{cryptography, security, encryption, authentication, hash functions, TLS, protocols, provable security}
\institute{University of Edinburgh, School of Informatics}
\instituteimage{images/edu-white.png}
\date{\today}
\coversubtitle{INFR10067\\Fall 2025}
\coverpartname{Secure communications}
\covertopicname{Public Key Infrastructure}
\coverwebsite{}

\mode<presentation>
{
  \usetheme{default}
  \usecolortheme{default}
  \usefonttheme{default}
  \setbeamertemplate{navigation symbols}{}
  \setbeamertemplate{caption}[empty]
  \setbeamertemplate{footline}[frame number]
}


\iffalse
\setbeamertemplate{blocks}[rounded][shadow=true]
\setbeamercolor{block title}{fg=black, bg=Grey!50}
\setbeamercolor{block body}{fg=black, bg=Grey!25}
\setbeamercolor{block title alerted}{fg=black, bg=Red!20}
\setbeamercolor{block body alerted}{fg=black, bg=Red!15}
\setbeamercolor{block title example}{fg=black, bg=SeaGreen!20}
\setbeamercolor{block body example}{fg=black, bg=Green!15}
\fi

%\usepackage[english]{babel}
\usepackage[utf8x]{inputenc}

\usepackage[lambda,advantage,operators,sets,adversary,landau,probability,notions,logic,ff,mm,primitives,events,complexity,asymptotics,keys]{cryptocode}

\usepackage{graphicx}
\usepackage{tikz}
\usetikzlibrary{arrows.meta}
\usetikzlibrary{shapes.geometric}
\pgfdeclarelayer{background}
\pgfsetlayers{background,main}

\usepackage{scrextend} 

\usepackage{fancybox}

\usepackage{pgfpages}

\usepackage{caption}
\usepackage{subcaption}

\usepackage{mdframed}
\usepackage{xspace}
\newcommand\Fontvi{\fontsize{5}{5.2}\selectfont}
\definecolor{ao}{rgb}{0.0, 0.6, 0.0}
\definecolor{bleu}{rgb}{0.2,0.2,0.7}

\newcommand{\Inred}[1]{\textcolor{BrickRed}{#1}}
\newcommand{\Inblue}[1]{\textcolor{bleu}{#1}}
\newcommand{\Ingreen}[1]{\textcolor{SeaGreen}{#1}}

\newcommand{\V}[1]{\texttt{V\textsubscript{{#1}}}\xspace}
\newcommand{\Vs}[1]{$\texttt{V}'_\texttt{{#1}}$\xspace}
\newcommand{\BB}{\texttt{BB}\xspace}
\newcommand{\EA}{\texttt{EA}\xspace}
\newcommand{\A}{\adv\xspace}
\newcommand{\F}[1]{$\mathcal{F}_\textsc{{#1}}$\xspace}
\renewcommand{\S}{\sdv\xspace}
\newcommand{\E}{$\mathcal{E}$\xspace}
\renewcommand{\P}{$P$\xspace}
\newcommand{\D}{\ddv\xspace}

\newcommand{\Pt}[2]{$\Pi^{{#2}}_{\textsc{{#1}}}$\xspace}

\newcommand{\Enc}{\texttt{Enc}}
\newcommand{\CheckShares}{\texttt{Check}}

\newcommand{\ASetup}{\texttt{ASetup}}
\newcommand{\AEval}{\texttt{AEval}}
\newcommand{\AWit}{\texttt{AWit}}
\newcommand{\AVer}{\texttt{AVer}}

\newcommand{\CSetup}{\texttt{CSetup}}
\newcommand{\CCommit}{\texttt{CCommit}}
\newcommand{\COpen}{\texttt{COpen}}

\newcommand{\SGen}{\texttt{SGen}}
\newcommand{\SSign}{\texttt{SSign}}
\newcommand{\SVer}{\texttt{SVer}}

\newcommand{\SoKSetup}{\texttt{SoKSetup}}
\newcommand{\SoKSign}{\texttt{SoKSign}}
\newcommand{\SoKVer}{\texttt{SoKVer}}

\newcommand{\TLEnc}{\texttt{TLEnc}}
\newcommand{\TLDec}{\texttt{TLDec}}

\newcommand{\frameindent}{\hspace*{\baselineskip}}

\colorlet{gris}{black!50!white}
\def\engris#1{\textcolor{gris}{#1}}

\colorlet{jaune}{yellow!90!black}
\def\enjaune#1{\textcolor{jaune!70!black}{#1}}

\colorlet{vert}{green!45!black}
\def\envert#1{\textcolor{vert}{#1}}

\colorlet{bleu}{blue!70!black}
\def\enbleu#1{\textcolor{bleu}{#1}}

\colorlet{rouge}{red!75!black}
\def\enrouge#1{\textcolor{rouge}{#1}}

\newcommand{\rouge}{\only{\color{rouge}}}

\xdefinecolor{violet}{rgb}{.5,0,.7}
\def\enviolet#1{\textcolor{violet}{#1}}

\xdefinecolor{violetft}{rgb}{0.199,0., 0.398}
\def\envioletft#1{\textcolor{violetft}{#1}}

\xdefinecolor{monbleu}{rgb}{0,.2,.54}
\newcommand{\monbleu}{\only{\color{monbleu}}}

\newcommand{\Ccal}{\mathcal{C}}
\newcommand{\Kcal}{\mathcal{K}}
\newcommand{\Mcal}{\mathcal{M}}
\newcommand{\Tcal}{\mathcal{T}}
\newcommand{\Xcal}{\mathcal{X}}

\usepackage{wasysym}
\usepackage{amsmath}
\usepackage{amssymb}
\usepackage{amsfonts}
\usepackage{ragged2e}
%\usepackage{msc}
\usepackage{color}
\usepackage{stmaryrd}
\usepackage{eurosym}



\title{Cryptography: public key infrastructure}

\author{Markulf Kohlweiss \\
  School of Informatics \\
  University of Edinburgh}



\date{}

\begin{document}

\begin{frame}
  \maketitle
\end{frame}

\begin{frame}
  \frametitle{Public keys}

  \begin{figure}[c]
    \centering
    \includegraphics[scale=0.2]{Images/asyouwish1.pdf}
    \caption{How does Alice trust that $pk_{Amazon}$ is Amazon's public key?}
  \end{figure}
  \bigskip

  \begin{block}{}
    Public-key encryption schemes are secure only if the authenticity of the public key is assured
  \end{block}
\end{frame}

\begin{frame}{Distribution of public keys}

  \begin{enumerate}
  \item \enbleu{Public announcements -} participants broadcast their public key

    \enrouge{\Large{\frownie}\ \normalsize does not defend against forgeries}
    \medskip{}

  \item \enbleu{Publicly available directories -} participants publish their public key on public directories

    \enrouge{\Large{\frownie}\ \normalsize does not defend against forgeries}
    \medskip{}

  \item \enbleu{Public-key authority -} participants contact the authority for each public key it needs 

    \enrouge{\Large{\frownie}\ \normalsize bottleneck in the system}
    \medskip{}

  \item \enbleu{public-key certificates -} CAs issue certificates to participants on their public key 

    \envert{\Large{\smiley}\ \normalsize as reliable as public-key authority but avoiding the bottleneck}
    \medskip{}

  \end{enumerate}
\end{frame}

\begin{frame}
  \frametitle{Public key certificates}
  \vspace{-0.65cm}

  \begin{figure}[c]
    \centering
    \includegraphics[scale=0.3]{Images/PublicKeyCertificate.pdf}
    \caption{\scriptsize image from Cryptography and Network Security - Principles and Practice - William Stallings }
  \end{figure}
  \vspace{-0.25cm}


  \begin{block}{A certificate consists mainly of }
    \begin{itemize}
    \item a \enbleu{public key}
    \item a \enbleu{subject} identifying the owner of the key
    \item a \enbleu{signature} by the CA on the key and the subject binding them together

      \enrouge{the CA is trusted}
    \end{itemize}
  \end{block}

\end{frame}

\begin{frame}
  \frametitle{X.509 certificates}

  \begin{figure}[c]
    \centering
    \includegraphics[scale=0.3]{Images/X509certificates.pdf}
    \caption{\scriptsize image from Cryptography and Network Security - Principles and Practice - William Stallings}
  \end{figure}

  \begin{itemize}
  \item X.509 defines a framework for the provision of authentication services
  \item Used by many applications such as TLS
  \end{itemize}
\end{frame}

\begin{frame}
  \frametitle{Public key certificates}
  \begin{figure}[c]
    \centering
    \includegraphics[scale=0.2]{Images/asyouwish2.pdf}
    \caption{Alice can now verify Amazon's certificate}
  \end{figure}

\end{frame}

\begin{frame}
  \frametitle{Using public key certificates to secure the Internet}
  \vspace{-0.45cm}
  \begin{center}
    \includegraphics[scale=0.45]{Images/AmazonWebsite.pdf}
  \end{center}


  \begin{block}{A very important implicit assumption}
    The browser is trusted to be ``secure''
  \end{block}
\end{frame}

\begin{frame}
  \frametitle{Amazon's certificate}
    \vspace{-0.45cm}
  \begin{center}
    \includegraphics[scale=0.3]{Images/AmazonCertificate.pdf}
  \end{center}

\end{frame}

\begin{frame}{Browser root certificates}
      \vspace{-0.45cm}
  \begin{center}
    \includegraphics[scale=0.3]{Images/RootCertificates.pdf}
  \end{center}

\end{frame}

\begin{frame}{Chain of trust}
    \vspace{-0.45cm}
  \begin{figure}[c]
    \centering
    \includegraphics[scale=0.5]{Images/Hierarchy.pdf}
    \caption{\scriptsize X.509 Hierarchy - image from Cryptography and Network Security - Principles and Practice - William Stallings}
  \end{figure}

  \begin{block}{}
    \begin{itemize}
    \item Having a single CA sign all certificates is not practical
    \item Instead a root CA signs certificates for level 1 CAs, level
      1 CAs sign certificates for level 2 CAs, \emph{etc}
    \end{itemize}
  \end{block}

\end{frame}

\begin{frame}{Self-signed certificates}

  \begin{center}
    \includegraphics[scale=0.4]{Images/selfsignedcertificate.pdf}
  \end{center}

\end{frame}

\begin{frame}
      \vspace{-0.45cm}
  \frametitle{The Lenovo Superfish scandal (February 2015)}
  \begin{center}
    \includegraphics[scale=0.65]{Images/LenovoSuperfish.pdf}
  \end{center}
\end{frame}

\begin{frame}
  \frametitle{And another scandal (September 2016)}

  \begin{center}
    \includegraphics[scale=0.85]{Images/WoSign.pdf}
  \end{center}

\end{frame}

\begin{frame}{Revocation}
    \vspace{-0.45cm}
    \begin{itemize}
    \item A certificate needs to be revoked if the corresponding private key has been compromised.\bigskip
      
    \item   Certificate Revocation Lists (CRLs) are the solution adopted in X.509.\bigskip
      
    \item Online Certificate Status Protocol (OCSP) stapling is the modern solution to this problem. 
    \end{itemize}

  \end{frame}

\end{document}
